\section{High centers and pointed factor rigidity}
\label{sec:factors}

The exact skeleton furnishes centers at which every bounded off-center
coordinate has stabilized.  This is the step that collapses arbitrary
finite context to a letter map.

\begin{lemma}[High-center identity]
\label{lem:high-center}
  Let $c_n=r_n$.  For every nonzero $j\in\Z$ and every
  $n>\nu_p(j)$,
  \begin{equation}
    \nu_p\bigl(L_p(c_n+j)\bigr)=\nu_p(j).
    \label{eq:high-center}
  \end{equation}
\end{lemma}

\begin{proof}
  Put $e=\nu_p(j)$ and write $j=p^e a$, where $p\nmid a$; the integer $a$
  may be negative.  By \cref{eq:center-identities},
  \begin{equation}
    L_p(c_n+j)=p^n+(p-1)j
      =p^e\bigl[p^{n-e}+(p-1)a\bigr].
    \label{eq:high-center-factor}
  \end{equation}
  Since $n-e\geq1$, the bracket is congruent to
  $(p-1)a\equiv-a$ modulo $p$.  It is not divisible by $p$, proving
  \cref{eq:high-center}.  The calculation uses no inverse and applies to
  prime and composite bases alike.
\end{proof}

We use the Curtis--Hedlund--Lyndon theorem in the following standard form:
a continuous shift-commuting map between subshifts is induced by a finite
local rule \citep{Hedlund1969}.  Padding its window if necessary, we may
take a symmetric radius $R$.

\begin{theorem}[Arbitrary-radius pointed factors collapse to letters]
\label{thm:factor-collapse}
  Let $u$ and $v$ be frozen directives at the same base $p$.  A pointed
  factor $F:T_{p,u}\to T_{p,v}$ exists if and only if there is a
  surjective letter map $\lambda:\calA\to\calB$ satisfying
  $v_n=\lambda(u_n)$ for all $n\in\Z$.  When it exists, $\lambda$ is
  unique and
  \begin{equation}
    F(z)(k)=\lambda(z(k))
    \qquad(z\in X_{p,u},\ k\in\Z).
    \label{eq:factor-coordinatewise}
  \end{equation}
  In particular, the pointed factor itself is unique.
\end{theorem}

\begin{proof}
  Suppose first that $F$ is a pointed factor.  Choose $R\geq0$ and a local
  rule $\phi$ such that
  \begin{equation}
    F(z)(k)=\phi\bigl(z[k-R,k+R]\bigr).
    \label{eq:chl-rule}
  \end{equation}
  If $R>0$, set
  \begin{equation}
    M=\max\{\nu_p(j):0<|j|\leq R\};
    \label{eq:window-threshold}
  \end{equation}
  if $R=0$, put $M=-1$.  For every $n>M$, the center coordinate is
  $x_{p,u}(c_n)=u_n$, while \cref{lem:high-center} gives
  \begin{equation}
    x_{p,u}(c_n+j)=u_{\nu_p(j)}
    \qquad(0<|j|\leq R).
    \label{eq:high-window-normal-form}
  \end{equation}
  Thus all entries in the radius-$R$ window except its center are
  independent of $n$.  Every source letter occurs at arbitrarily large
  indices because $u$ is periodic with exact support.  For $a\in\calA$,
  let $W_a$ be the stabilized window with center $a$ and define
  \begin{equation}
    \lambda(a)=\phi(W_a).
    \label{eq:lambda-definition}
  \end{equation}
  This is independent of the chosen large occurrence of $a$.

  Pointedness and $L_p(c_n)=p^n$ now yield, for $n>M$,
  \begin{equation}
    v_n=x_{p,v}(c_n)=F(x_{p,u})(c_n)=\lambda(u_n).
    \label{eq:large-directive-relation}
  \end{equation}
  Let $h_u$ and $h_v$ be the least directive periods and set
  $H=\lcm(h_u,h_v)$.  Given any $n\in\Z$, choose $s$ so large that
  $n+sH>M$.  Periodicity extends \cref{eq:large-directive-relation} to
  \begin{equation}
    v_n=v_{n+sH}=\lambda(u_{n+sH})=\lambda(u_n).
    \label{eq:all-directive-relation}
  \end{equation}
  Therefore, at every $k\in\Z$,
  \begin{align}
    x_{p,v}(k)
      &=v_{\nu_p(L_p(k))}
        =\lambda\bigl(u_{\nu_p(L_p(k))}\bigr)
        =\lambda(x_{p,u}(k)).
    \label{eq:basepoint-coordinate-map}
  \end{align}
  The original map $F$ and $\lambda^{\Z}$ agree on the distinguished
  point, hence on its shift orbit by equivariance, and hence on its closure
  by continuity.  This proves \cref{eq:factor-coordinatewise}.

  Every target letter is some $v_n$ by exact target support.
  \Cref{eq:all-directive-relation} therefore makes $\lambda$ surjective.
  Exact source support also makes it unique: if $u_n=a$, then necessarily
  $\lambda(a)=v_n$.  Consequently $F$ is unique as well.

  Conversely, suppose that a surjective letter map satisfies
  $v=\lambda\circ u$.  Its coordinatewise extension is continuous and
  shift commuting and sends $x_{p,u}$ to $x_{p,v}$.  We verify its image
  by two inclusions.  If $y\in X_{p,u}$, choose integers $n_i$ such that
  $\sigma^{n_i}x_{p,u}\to y$.  Continuity and equivariance give
  \[
    \lambda^{\Z}(y)
      =\lim_i\sigma^{n_i}x_{p,v}\in X_{p,v},
  \]
  so $\lambda^{\Z}(X_{p,u})\subseteq X_{p,v}$.  Conversely,
  $\lambda^{\Z}(X_{p,u})$ is compact and therefore closed, and it contains
  $\Orb(x_{p,v})$ because it contains $x_{p,v}$ and is shift invariant.
  It consequently contains $\cl\,\Orb(x_{p,v})=X_{p,v}$.  The two
  inclusions prove equality, so the coordinatewise extension is onto and
  is the required pointed factor.
\end{proof}

\begin{corollary}[Pointed conjugacy]
\label{cor:pointed-conjugacy}
  Two systems $T_{p,u}$ and $T_{p,v}$ are pointedly conjugate exactly when
  their directives differ by a bijective relabeling of letters.  The
  conjugacy has no phase shift.
\end{corollary}

\begin{proof}
  A bijective letter map has the inverse coordinate map.  Conversely, the
  inverse of a pointed conjugacy is again pointed and same-base.  Apply
  \cref{thm:factor-collapse} in both directions.  The two resulting letter
  maps compose to the identity on every directive letter, because every
  such letter occurs at a center $c_n$.  They are therefore bijections.
\end{proof}

\begin{remark}[Where the argument uses the scope]
\label{rem:factor-scope}
  Same-base scope provides common centers $c_n$ in
  \cref{eq:large-directive-relation}; pointedness identifies the local-rule
  output there with $v_n$ rather than a shifted target coordinate; and an
  in-family target supplies its directive and exact support.  Removing any
  of these hypotheses breaks a displayed step.  The theorem does not assert
  that wrong-base or nonpointed maps are impossible in broader categories.
\end{remark}
